\section{Emergence of hybrid cloud}
\label{sec:Emergence}

In recent years, hybrid cloud has been widely adopted, transforming operations in many enterprise IT. As a result, this transformation helps organizations to be capable of integrating private and public cloud infrastructures into a single computing environment \cite{Mell2011NIST}, \cite{Liu2011NISTRA}. This deployment model allows enterprises to optimize resource utilization, improve business continuity, and maintain diverse regulatory and operational requirements while ensuring the flexibility to successfully deploy workloads across heterogeneous infrastructures. Moreover, \cite{Liu2011NISTRA} defines a base hybrid cloud that facilitates portability and interoperability of the application between multiple cloud environments. In this way, organizations can benefit from the scalability and elasticity of public cloud services without giving up control over sensitive data and mission-critical applications.

However, the integration of multiple cloud environments significantly increases the complexity of infrastructure management \cite{Armbrust2010Cloud}, \cite{Flexera2025StateCloud}, \cite{Liu2011NISTRA}. In general, hybrid cloud models consist of heterogeneous technologies, distributed resources, and multiple security models that must operate as a cohesive system. Consequently, organizations are required to maintain consistent governance, identity and access management, and security policies in both private and public cloud in hybrid cloud. As cloud-native applications and distributed services continue to scale, ensuring operational consistency and maintaining a security standard become increasingly difficult \cite{HashiCorp2025StateCloud}.

Furthermore, hybrid cloud environments require strict coordination between infrastructure management, application deployment, and security operations \cite{Armbrust2010Cloud}, \cite{HashiCorp2025StateCloud}. Traditional security approaches, which are based on periodic assessments or isolated security modules, are incapable of addressing the dynamic nature of hybrid cloud infrastructures. Instead, security controls must be integrated throughout the system lifecycle to support continuous monitoring, automated vulnerability assessment, and policy enforcement. Consequently, organizations require operational processes that are integrated with automation, infrastructure management, and continuous security practices to effectively manage hybrid cloud environments while maintaining security, and compliance \cite{Flexera2025StateCloud}.